
\documentclass{article}
\usepackage{colortbl}
\usepackage{makecell}
\usepackage{multirow}
\usepackage{supertabular}

\begin{document}

\newcounter{utterance}

\centering \large Interaction Transcript for game `cladder', experiment `full\_v1.5\_default', episode 7754 with qwen3.5{-}27b.
\vspace{24pt}

{ \footnotesize  \setcounter{utterance}{1}
\setlength{\tabcolsep}{0pt}
\begin{supertabular}{c@{$\;$}|p{.15\linewidth}@{}p{.15\linewidth}p{.15\linewidth}p{.15\linewidth}p{.15\linewidth}p{.15\linewidth}}
    \# & \multicolumn{2}{c}{Player} && \multicolumn{2}{c}{Game Master} \\
    \hline

    \theutterance \stepcounter{utterance}  
    & & & \multicolumn{4}{p{0.6\linewidth}}{
        \cellcolor[rgb]{0.9,0.9,0.9}{
            \makecell[{{p{\linewidth}}}]{
                \texttt{\tiny{[P1$\langle$GM]}}
                \texttt{You are an expert in causal inference. The following question is not a typical commonsense query, but rather a meticulously designed question created by a professor specializing in causal inference, intended to assess the students' mastery of the course content.} \\
\texttt{Imagine a self{-}contained, hypothetical world with only the following conditions, and without any unmentioned factors or causal relationships: Zory has a direct effect on zuph and glimx. Zuph has a direct effect on jyka. Jyka has a direct effect on glimx. Zory is unobserved. The overall probability of zuph is 8\%. The probability of not zuph and glimx is 58\%. The probability of zuph and glimx is 1\%. Is the chance of glimx smaller when observing zuph?} \\
\texttt{Start your answer with 'Yes' or 'No', followed by additional reasoning or evidence to support your explanation.} \\
            }
        }
    }
     \\ \\

    \theutterance \stepcounter{utterance}  
    & \multicolumn{4}{p{0.6\linewidth}}{
        \cellcolor[rgb]{0.561,0.737,0.561}{
            \makecell[{{p{\linewidth}}}]{
                \texttt{\tiny{[P1$\rangle$GM]}}
                \texttt{No} \\
\\ 
\texttt{To determine whether the chance of *glimx* is smaller when observing *zuph*, we must calculate and compare the conditional probabilities \$P(\textbackslash{}text\{glimx\} \textbackslash{}mid \textbackslash{}text\{zuph\})\$ and \$P(\textbackslash{}text\{glimx\} \textbackslash{}mid \textbackslash{}neg\textbackslash{}text\{zuph\})\$.} \\
\\ 
\texttt{**1. Calculate \$P(\textbackslash{}text\{glimx\} \textbackslash{}mid \textbackslash{}text\{zuph\})\$:**} \\
\texttt{We are given:} \\
\texttt{*   \$P(\textbackslash{}text\{zuph\}) = 0.08\$} \\
\texttt{*   \$P(\textbackslash{}text\{zuph\} \textbackslash{}cap \textbackslash{}text\{glimx\}) = 0.01\$} \\
\\ 
\texttt{Using the definition of conditional probability:} \\
\texttt{\$\$P(\textbackslash{}text\{glimx\} \textbackslash{}mid \textbackslash{}text\{zuph\}) = \textbackslash{}frac\{P(\textbackslash{}text\{zuph\} \textbackslash{}cap \textbackslash{}text\{glimx\})\}\{P(\textbackslash{}text\{zuph\})\} = \textbackslash{}frac\{0.01\}\{0.08\} = 0.125\$\$} \\
\\ 
\texttt{**2. Calculate \$P(\textbackslash{}text\{glimx\} \textbackslash{}mid \textbackslash{}neg\textbackslash{}text\{zuph\})\$:**} \\
\texttt{We are given:} \\
\texttt{*   \$P(\textbackslash{}neg\textbackslash{}text\{zuph\} \textbackslash{}cap \textbackslash{}text\{glimx\}) = 0.58\$} \\
\\ 
\texttt{First, find the probability of not *zuph*:} \\
\texttt{\$\$P(\textbackslash{}neg\textbackslash{}text\{zuph\}) = 1 {-} P(\textbackslash{}text\{zuph\}) = 1 {-} 0.08 = 0.92\$\$} \\
\\ 
\texttt{Now, calculate the conditional probability:} \\
\texttt{\$\$P(\textbackslash{}text\{glimx\} \textbackslash{}mid \textbackslash{}neg\textbackslash{}text\{zuph\}) = \textbackslash{}frac\{P(\textbackslash{}neg\textbackslash{}text\{zuph\} \textbackslash{}cap \textbackslash{}text\{glimx\})\}\{P(\textbackslash{}neg\textbackslash{}text\{zuph\})\} = \textbackslash{}frac\{0.58\}\{0.92\} \textbackslash{}approx 0.630\$\$} \\
\\ 
\texttt{**3. Comparison:**} \\
\texttt{*   \$P(\textbackslash{}text\{glimx\} \textbackslash{}mid \textbackslash{}text\{zuph\}) = 0.125\$} \\
\texttt{*   \$P(\textbackslash{}text\{glimx\} \textbackslash{}mid \textbackslash{}neg\textbackslash{}text\{zuph\}) \textbackslash{}approx 0.630\$} \\
\\ 
\texttt{Since \$0.125 < 0.630\$, the probability of *glimx* is indeed significantly smaller when *zuph* is observed compared to when it is not.} \\
\\ 
\texttt{**Causal Context Note:**} \\
\texttt{Although the question asks about the statistical association (which is negative here), the causal structure described (\$Zory \textbackslash{}to Zuph\$, \$Zory \textbackslash{}to Glimx\$, \$Zuph \textbackslash{}to Jyka \textbackslash{}to Glimx\$) creates a "collider" or "fork" scenario depending on the specific paths. Specifically, \$Zory\$ is a common cause of both \$Zuph\$ and \$Glimx\$. Since \$Zory\$ is unobserved, conditioning on \$Zuph\$ (a descendant of the common cause) induces a correlation. However, the direct calculation of the probabilities provided in the prompt is sufficient to answer the specific question about the magnitude of the chance, regardless of the underlying causal mechanism's directionality in this specific numerical instance. The data explicitly shows a negative association.} \\
            }
        }
    }
    & & \\ \\

    \theutterance \stepcounter{utterance}  
    & & & \multicolumn{2}{p{0.3\linewidth}}{
        \cellcolor[rgb]{0.9,0.9,0.9}{
            \makecell[{{p{\linewidth}}}]{
                \texttt{\tiny{[GM$|$GM]}}
                \texttt{No} \\
            }
        }
    }
    & & \\ \\

    \theutterance \stepcounter{utterance}  
    & & & \multicolumn{2}{p{0.3\linewidth}}{
        \cellcolor[rgb]{0.9,0.9,0.9}{
            \makecell[{{p{\linewidth}}}]{
                \texttt{\tiny{[GM$|$GM]}}
                \texttt{yes} \\
            }
        }
    }
    & & \\ \\

    \theutterance \stepcounter{utterance}  
    & & & \multicolumn{2}{p{0.3\linewidth}}{
        \cellcolor[rgb]{0.9,0.9,0.9}{
            \makecell[{{p{\linewidth}}}]{
                \texttt{\tiny{[GM$|$GM]}}
                \texttt{game\_result = LOSE} \\
            }
        }
    }
    & & \\ \\

\end{supertabular}
}

\end{document}
